% Options for packages loaded elsewhere
\PassOptionsToPackage{unicode}{hyperref}
\PassOptionsToPackage{hyphens}{url}
\documentclass[
  12pt,
]{ctexart}
\usepackage{xcolor}
\usepackage{amsmath,amssymb}
\setcounter{secnumdepth}{-\maxdimen} % remove section numbering
\usepackage{iftex}
\ifPDFTeX
  \usepackage[T1]{fontenc}
  \usepackage[utf8]{inputenc}
  \usepackage{textcomp} % provide euro and other symbols
\else % if luatex or xetex
  \usepackage{unicode-math} % this also loads fontspec
  \defaultfontfeatures{Scale=MatchLowercase}
  \defaultfontfeatures[\rmfamily]{Ligatures=TeX,Scale=1}
\fi
\usepackage{lmodern}
\ifPDFTeX\else
  % xetex/luatex font selection
\fi
% Use upquote if available, for straight quotes in verbatim environments
\IfFileExists{upquote.sty}{\usepackage{upquote}}{}
\IfFileExists{microtype.sty}{% use microtype if available
  \usepackage[]{microtype}
  \UseMicrotypeSet[protrusion]{basicmath} % disable protrusion for tt fonts
}{}
\makeatletter
\@ifundefined{KOMAClassName}{% if non-KOMA class
  \IfFileExists{parskip.sty}{%
    \usepackage{parskip}
  }{% else
    \setlength{\parindent}{0pt}
    \setlength{\parskip}{6pt plus 2pt minus 1pt}}
}{% if KOMA class
  \KOMAoptions{parskip=half}}
\makeatother
\usepackage{longtable,booktabs,array}
\usepackage{calc} % for calculating minipage widths
% Correct order of tables after \paragraph or \subparagraph
\usepackage{etoolbox}
\makeatletter
\patchcmd\longtable{\par}{\if@noskipsec\mbox{}\fi\par}{}{}
\makeatother
% Allow footnotes in longtable head/foot
\IfFileExists{footnotehyper.sty}{\usepackage{footnotehyper}}{\usepackage{footnote}}
\makesavenoteenv{longtable}
\setlength{\emergencystretch}{3em} % prevent overfull lines
\providecommand{\tightlist}{%
  \setlength{\itemsep}{0pt}\setlength{\parskip}{0pt}}
\usepackage{bookmark}
\IfFileExists{xurl.sty}{\usepackage{xurl}}{} % add URL line breaks if available
\urlstyle{same}
\hypersetup{
  hidelinks,
  pdfcreator={LaTeX via pandoc}}

\author{SCX}
\date{2026}

\begin{document}

\section{START\_HERE\_CODEX --- SCX 项目 AI
恢复入口}\label{start_here_codex-scx-ux9879ux76ee-ai-ux6062ux590dux5165ux53e3}

\begin{quote}
最后更新：2026-06-27

\textbf{如果你是 AI（Claude Code / GPT /
其他）被用户拉入这个项目，这是你的第一份阅读材料。}
读完这份文件再行动。不要跳过。
\end{quote}

\begin{center}\rule{0.5\linewidth}{0.5pt}\end{center}

\subsection{项目根目录}\label{ux9879ux76eeux6839ux76eeux5f55}

\texttt{G:\textbackslash{}Xiaogan\_Supercomputing\_data\textbackslash{}SCX}

\subsection{一句话定位}\label{ux4e00ux53e5ux8bddux5b9aux4f4d}

\textbf{SCX: State-Conditioned eXpertise} ---
状态条件专家性框架。核心主张：专家可靠性不是全局常数，而是状态条件函数。数据价值不由样本固有属性决定，而由状态、专家可靠性和当前模型缺陷共同决定。

\begin{center}\rule{0.5\linewidth}{0.5pt}\end{center}

\subsection{版本状态}\label{ux7248ux672cux72b6ux6001}

\textbf{v0.4.0-pre} (2026-06-27)

\begin{longtable}[]{@{}ll@{}}
\toprule\noalign{}
指标 & 数值 \\
\midrule\noalign{}
\endhead
\bottomrule\noalign{}
\endlastfoot
Python 文件 & 35 (\texttt{src/scx/} 6 子包, 34 模块) \\
单元测试 & \textbf{370 tests}（全部通过） \\
核心定理 & \textbf{2}：定理 1（噪声检测保证）+ 定理
2（弱特征失效下界） \\
旧命题 & 6 → 重构中（V(s) 已标记 deprecated） \\
GPU & ❌ 无（torch CPU-only，等显卡） \\
超算 & ✅ cxshu（仅 EGP DFT 计算用） \\
\end{longtable}

\begin{center}\rule{0.5\linewidth}{0.5pt}\end{center}

\subsection{用户习惯与偏好}\label{ux7528ux6237ux4e60ux60efux4e0eux504fux597d}

\subsubsection{沟通风格}\label{ux6c9fux901aux98ceux683c}

\begin{itemize}
\tightlist
\item
  \textbf{中文为主}，技术术语可中英混用
\item
  偏好\textbf{独立思考}：不要盲从 GPT
  的建议，要从理论和审稿人视角出发独立判断
\item
  重视\textbf{严密性}：理论证明 \textgreater{} 实验 \textgreater{}
  叙事。反对''空中楼阁''
\item
  喜欢用 \texttt{{[}{[}wikilink{]}{]}} 和 Obsidian 图谱组织知识
\end{itemize}

\subsubsection{工作方式}\label{ux5de5ux4f5cux65b9ux5f0f}

\begin{itemize}
\tightlist
\item
  \textbf{多 Agent 并行}：能用 agent 并发的事不要串行
\item
  \textbf{先规划后执行}：复杂任务先写 plan → 用户审批 → 再动手
\item
  \textbf{文件合并癖}：多个文件讨论同一主题 → 合并为单一权威来源
\item
  \textbf{TODO 狂热}：喜欢细粒度、可追踪、分优先级的任务清单
\item
  \textbf{里程碑驱动}：喜欢看到进度（M1→M2→M3\ldots）
\end{itemize}

\subsubsection{价值观}\label{ux4ef7ux503cux89c2}

\begin{itemize}
\tightlist
\item
  \textbf{学术优先权 \textgreater{} 商业}：先发论文占坑，再考虑商业
\item
  \textbf{但商业不是不重要}：论文公开 L1-L2（数学+算法），工程 L4-L5
  闭源保留
\item
  \textbf{诚实面对失败}：弱特征失效不是 bug
  是定理。审稿人喜欢看到作者知道方法的边界
\item
  \textbf{Nature 是目标不是梦}：冲击正刊需要''惊人发现 + 理论保证 +
  多领域证据''
\end{itemize}

\subsubsection{当前心态}\label{ux5f53ux524dux5fc3ux6001}

\begin{itemize}
\tightlist
\item
  等显卡焦虑但不想浪费时间 → 把能做的理论/写作/CPU实验先做了
\item
  项目从 EGP 到 SCX 迭代太快，需要 Obsidian 库来追溯思想演化
\item
  Paper 2（EGP gauge fixing）是最接近投稿的，想赶紧投出去建立根据地
\end{itemize}

\begin{center}\rule{0.5\linewidth}{0.5pt}\end{center}

\subsection{5 篇论文谱系}\label{ux7bc7ux8bbaux6587ux8c31ux7cfb}

\begin{quote}
唯一权威来源：\texttt{knowledge/05\_决策/论文规划.md}
\end{quote}

\begin{verbatim}
Paper 2 (先发, npj)    Paper 1 (旗舰, Nature)       Paper 3 (TMLR)    Paper 4 (JMLR)   Paper 5 (远期)
Gauge Fixing           Data Quality > Architecture   噪声检测理论        压缩保真理论       跨领域深挖
    ↓                        ↓                          ↓                  ↓               ↓
 纯应用                   Thm 1+2 轻量版              Thm 1+2 放松版     Thm 3 完善版      合作方驱动
 已有草稿+8图             跨领域 4 domain             5→3 假设           coreset 对话       Sim/Health
 1-2周可投稿              方法完整框架                 Bernstein/Talagrand Feldman-Langberg  /Robot
\end{verbatim}

\begin{longtable}[]{@{}
  >{\raggedright\arraybackslash}p{(\linewidth - 8\tabcolsep) * \real{0.0882}}
  >{\raggedright\arraybackslash}p{(\linewidth - 8\tabcolsep) * \real{0.1765}}
  >{\raggedright\arraybackslash}p{(\linewidth - 8\tabcolsep) * \real{0.2647}}
  >{\raggedright\arraybackslash}p{(\linewidth - 8\tabcolsep) * \real{0.2941}}
  >{\raggedright\arraybackslash}p{(\linewidth - 8\tabcolsep) * \real{0.1765}}@{}}
\toprule\noalign{}
\begin{minipage}[b]{\linewidth}\raggedright
\#
\end{minipage} & \begin{minipage}[b]{\linewidth}\raggedright
论文
\end{minipage} & \begin{minipage}[b]{\linewidth}\raggedright
核心问题
\end{minipage} & \begin{minipage}[b]{\linewidth}\raggedright
目标期刊
\end{minipage} & \begin{minipage}[b]{\linewidth}\raggedright
状态
\end{minipage} \\
\midrule\noalign{}
\endhead
\bottomrule\noalign{}
\endlastfoot
\textbf{2} & Gauge-Normalized ACE for AlN & shared+correction ACE 的
gauge freedom 识别与 post-hoc 修复 & npj Comput. Mater. & ✅
已有草稿+8图，1-2周可投 \\
\textbf{1} & Data Quality \textgreater{} Model Architecture &
训练数据清洗 vs 架构改进：SCX 多专家一致性跨领域验证 & Nature Comput.
Sci. & 🔴 等 GPU 补齐实验 \\
\textbf{3} & Noise Detection Theory & Thm 1+2 放松假设版 + Dawid-Skene +
PAC-Bayes & TMLR / AISTATS & 📐 定理已证明 \\
\textbf{4} & Compression Fidelity Theory & Thm 3 (Prop 4 完善版) +
coreset theory & JMLR & 📐 需完善 bound \\
\textbf{5} & Cross-Domain Deep Dive & 选一个领域做 10× 深度验证 & 待定 &
❌ 远期 \\
\end{longtable}

\textbf{定理放置（方案 D）}：Paper 1 正文轻量陈述 + SI 完整证明。Paper
3+4 放松假设完整理论版。

\begin{center}\rule{0.5\linewidth}{0.5pt}\end{center}

\subsection{商业化策略}\label{ux5546ux4e1aux5316ux7b56ux7565}

\textbf{5 层资产模型}（来自与 GPT 的深度讨论）：

\begin{longtable}[]{@{}
  >{\raggedright\arraybackslash}p{(\linewidth - 6\tabcolsep) * \real{0.2500}}
  >{\raggedright\arraybackslash}p{(\linewidth - 6\tabcolsep) * \real{0.2500}}
  >{\centering\arraybackslash}p{(\linewidth - 6\tabcolsep) * \real{0.2500}}
  >{\raggedright\arraybackslash}p{(\linewidth - 6\tabcolsep) * \real{0.2500}}@{}}
\toprule\noalign{}
\begin{minipage}[b]{\linewidth}\raggedright
层级
\end{minipage} & \begin{minipage}[b]{\linewidth}\raggedright
内容
\end{minipage} & \begin{minipage}[b]{\linewidth}\centering
公开？
\end{minipage} & \begin{minipage}[b]{\linewidth}\raggedright
策略
\end{minipage} \\
\midrule\noalign{}
\endhead
\bottomrule\noalign{}
\endlastfoot
L1 数学理论 & Thm 1+2+3, 命题证明 & ✅ 全部发表 & 学术优先权 \\
L2 算法框架 & pseudocode, 核心流程 & ✅ 全部发表 & 学术优先权 \\
L3 最小代码 & reference implementation & ✅ 开源 & 复现论文用 \\
\textbf{L4 工程实现} & pipeline, dashboard, 云API & ❌ 闭源 &
\textbf{商业壁垒} \\
\textbf{L5 数据/专家库} & 校准参数, 训练好的势函数 & ❌ 闭源 &
\textbf{核心资产} \\
\end{longtable}

\textbf{操作顺序}：invention disclosure → 专利评估 → 提交专利申请 →
arXiv → 期刊投稿 → 开源最小代码 → 商业版

\textbf{关键认知}：公式定理全公开建立''SCX
是你发明的''学术事实。竞争对手可以读论文但不能复制工程调优和数据积累。Google
发表了 PageRank 公式，搜索引擎壁垒在爬虫和索引------不在公式。

\begin{center}\rule{0.5\linewidth}{0.5pt}\end{center}

\subsection{快速导航：Obsidian
知识库}\label{ux5febux901fux5bfcux822aobsidian-ux77e5ux8bc6ux5e93}

\begin{quote}
Obsidian vault 根目录：\texttt{knowledge/}

在 Obsidian 中打开 \texttt{knowledge/} 作为 vault，\texttt{Ctrl+G}
查看图谱。
\end{quote}

\begin{longtable}[]{@{}
  >{\raggedright\arraybackslash}p{(\linewidth - 2\tabcolsep) * \real{0.6000}}
  >{\raggedright\arraybackslash}p{(\linewidth - 2\tabcolsep) * \real{0.4000}}@{}}
\toprule\noalign{}
\begin{minipage}[b]{\linewidth}\raggedright
我想了解\ldots{}
\end{minipage} & \begin{minipage}[b]{\linewidth}\raggedright
去这里
\end{minipage} \\
\midrule\noalign{}
\endhead
\bottomrule\noalign{}
\endlastfoot
📘 \textbf{论文规划（唯一权威）} &
\texttt{knowledge/05\_决策/论文规划.md} \\
🧮 定理 1+2 完整证明 &
\texttt{theory/theorems/01\_noise\_detection\_guarantee.md},
\texttt{02\_weak\_feature\_failure.md} \\
📋 \textbf{任务清单（全面推进）} &
\texttt{knowledge/07\_任务/README.md} \\
🎯 里程碑追踪 & \texttt{knowledge/07\_任务/里程碑.md} \\
🔬 CPU 可做的实验 & \texttt{knowledge/07\_任务/NOW\_实验.md} \\
🧮 待证的理论 & \texttt{knowledge/07\_任务/NOW\_理论.md} \\
📝 论文推进任务 & \texttt{knowledge/07\_任务/NOW\_论文.md} \\
💻 代码工程任务 & \texttt{knowledge/07\_任务/NOW\_代码.md} \\
⏳ GPU 阻塞任务 & \texttt{knowledge/07\_任务/GPU\_等待.md} \\
🔍 代码审查报告 & \texttt{knowledge/07\_任务/代码审查\_定理同步.md} \\
📅 时间线 & \texttt{knowledge/01\_时间线/README.md} \\
🗺️ 思想演化地图 & \texttt{knowledge/00\_入口/思想演化地图.md} \\
📖 论文谱系 & \texttt{knowledge/00\_入口/论文谱系.md} \\
📊 项目状态 & \texttt{knowledge/03\_项目/SCX-框架.md} \\
💼 IP 策略 & \texttt{knowledge/05\_决策/IP策略.md} \\
🏠 知识库首页 & \texttt{knowledge/00\_入口/🏠\_Home.md} \\
\end{longtable}

\begin{center}\rule{0.5\linewidth}{0.5pt}\end{center}

\subsection{项目结构}\label{ux9879ux76eeux7ed3ux6784}

\begin{verbatim}
SCX/
├── knowledge/                # 📓 Obsidian 知识库（AI 的主要导航界面）
│   ├── .obsidian/            # Obsidian 配置
│   ├── 00_入口/              # MOC：首页、思想演化地图、论文谱系
│   ├── 01_时间线/            # 每日开发日志（6/14 → 6/27）
│   ├── 02_概念/              # 原子概念笔记（互相 [[链接]]）
│   ├── 03_项目/              # 项目状态面板
│   ├── 04_对话/              # GPT + Agent 讨论摘要
│   ├── 05_决策/              # 决策日志、IP策略、论文规划
│   ├── 06_收件箱/            # 临时笔记
│   └── 07_任务/              # 任务清单 + 里程碑
│
├── CodexKnowledge/           # AI 恢复入口 + 原文档案（本文件在此）
│   ├── START_HERE_CODEX.md  # ← 你正在读的文件
│   ├── SCX_核心定义.md       # 核心数学定义速查
│   ├── SCX_发展史与成就.md   # 发展历程（含 EGP-V2 agent 修订）
│   ├── SCX_思想扩展_综合方案.md
│   ├── SCX_架构清单.md       # 全项目文件清单
│   ├── SCX_TODO_论文框架.md  # 五阶段论文计划（五审查员产出）
│   ├── 决策日志.md           # 关键决策记录
│   ├── 工具状态.md           # 工具链状态
│   ├── 论文框架审查报告_五审查员综合.md
│   ├── 整理_IP保护与开源策略与商业模式.md
│   ├── 整理_LLM战略与发展路线图.md
│   ├── 整理_框架与模块与论文分布.md
│   ├── 与GPT的讨论2026-06-26-07-49.md  # 244KB GPT 对话（IP/商业/开源）
│   ├── 与gpt的对话202606261332.md      # 158KB GPT 对话（架构/多领域）
│   ├── agent_outputs/        # 8 份 Agent 分析文档
│   ├── archive/              # 过期文档
│   └── images/               # 讨论配图
│
├── theory/                   # 数学框架
│   ├── README.md             # 框架总览
│   ├── definitions/          # 6 个核心定义
│   ├── propositions/         # 6 个旧命题（重构中）
│   └── theorems/             # ✅ 2 个新定理（完整证明）
│       ├── 01_noise_detection_guarantee.md  # Thm 1: F1 ≥ 1 - exp(-2MΔ²)/η
│       └── 02_weak_feature_failure.md       # Thm 2: F1 ≤ baseline + √(2δ)
│
├── paper/                    # 论文相关
│   └── paper1_mlip/          # Paper 2 (EGP) LaTeX 草稿 + 8 张图 + supplementary
│
├── src/scx/                  # SCX Python 包（6 子包, 34 模块, ~16,000 行）
│   ├── core/                 # SCXFramework, Config, Metrics
│   ├── state/                # StateDiscovery, StateAssignment
│   ├── expert/               # ExpertRegistry, ExpertReliability, ExpertRouter, ExpertConflict
│   ├── valuation/            # DataClassifier, Learnability, Noise, Redundancy, StateValue, Influence
│   └── action/               # Acquisition, Compress, Policy
│
├── tests/                    # 370 tests (9 文件)
├── experiments/              # 验证实验
│   ├── synthetic/            # 合成 2D 实验
│   ├── cifar/                # CIFAR-10/100（SCX-Noise F1=0.617）
│   ├── mlip_case/            # AlN v3 两层分析（F1=0.585）
│   └── ml_benchmarks/        # 通用 ML 基准
├── scx-health/               # 医学数据估值子项目（MedMNIST, Routing 93.22%）
├── outputs/                  # 实验输出/图表
├── images/                   # 概念图
├── 势函数合并/               # 6 份早期研究想法文档
└── pyproject.toml
\end{verbatim}

\begin{center}\rule{0.5\linewidth}{0.5pt}\end{center}

\subsection{核心定理速查}\label{ux6838ux5fc3ux5b9aux7406ux901fux67e5}

\subsubsection{Theorem 1:
多专家一致性噪声检测保证}\label{theorem-1-ux591aux4e13ux5bb6ux4e00ux81f4ux6027ux566aux58f0ux68c0ux6d4bux4fddux8bc1}

\textbf{核心结果}：F1 ≥ 1 - (1/η) · Σ ρ\_s · exp(-2M · Δ\_s²)

\textbf{5 个假设}：(A1) 不相交训练集, (A2) 清洁数据条件独立, (A3)
有界损失, (A4) 均匀独立噪声, (A5) 状态同质性

\textbf{证明方法}：Hoeffding 不等式 + Chernoff 界 + 条件期望分解

\textbf{实践推论}：M ≥ 20 且平均专家错误率 ≤ 20\% → F1 ≥ 0.95

\textbf{代码位置}：\texttt{src/scx/valuation/state\_value.py} →
\texttt{noise\_detection\_f1\_bound()},
\texttt{noise\_consistency\_score()},
\texttt{optimal\_noise\_threshold()}

\subsubsection{Theorem 2:
弱特征失效下界}\label{theorem-2-ux5f31ux7279ux5f81ux5931ux6548ux4e0bux754c}

\textbf{核心结果}：F1 ≤ F1\_base + C\_F · √(2δ)，其中 δ = I(φ(X); s(X))

\textbf{证明方法}：Fano 不等式 + 耦合论证 + Pinsker 不等式 +
数据处理不等式

\textbf{实践推论}：特征互信息 δ → 0 时 SCX 退化到随机基线。解释了
DermaMNIST 上的失效

\textbf{代码位置}：\texttt{src/scx/valuation/state\_value.py} →
\texttt{feature\_strength\_diagnostic()}

\subsubsection{V(s) 状态}\label{vs-ux72b6ux6001}

旧 V(s) = r̄(s)·ρ(s)·L(s)·{[}1-D(s){]}·maxSCX(s) 已标记
\textbf{deprecated}（循环定义不可修复）。Thm 1+2
替代了其理论功能。代码中 \texttt{acquisition\_value()} 和
\texttt{compression\_value()} 仍可用但 emit DeprecationWarning。

\begin{center}\rule{0.5\linewidth}{0.5pt}\end{center}

\subsection{环境}\label{ux73afux5883}

\begin{longtable}[]{@{}
  >{\raggedright\arraybackslash}p{(\linewidth - 2\tabcolsep) * \real{0.5455}}
  >{\raggedright\arraybackslash}p{(\linewidth - 2\tabcolsep) * \real{0.4545}}@{}}
\toprule\noalign{}
\begin{minipage}[b]{\linewidth}\raggedright
资源
\end{minipage} & \begin{minipage}[b]{\linewidth}\raggedright
值
\end{minipage} \\
\midrule\noalign{}
\endhead
\bottomrule\noalign{}
\endlastfoot
Python &
\texttt{C:\textbackslash{}Users\textbackslash{}admin\textbackslash{}AppData\textbackslash{}Local\textbackslash{}hermes\textbackslash{}hermes-agent\textbackslash{}venv\textbackslash{}Scripts\textbackslash{}python.exe}
(3.11.9) \\
EGP venv (含 torch CPU) &
\texttt{D:/SHEprogram/EGP/.venv/Scripts/python.exe} \\
GPU & ❌ 无（torch CPU-only） \\
CPU ML & ✅ numpy, scipy, sklearn, pandas \\
SCX 安装 & \texttt{pip\ install\ -e\ .} (370 tests pass) \\
超算 & \texttt{ssh\ -p\ 24678\ cxshu@219.139.78.251} (AlN v3 DFT
已生成) \\
Obsidian vault & \texttt{knowledge/} \\
LaTeX & 无本地编译器，论文用 Overleaf \\
\end{longtable}

\begin{center}\rule{0.5\linewidth}{0.5pt}\end{center}

\subsection{AI 操作协议}\label{ai-ux64cdux4f5cux534fux8bae}

\subsubsection{每次会话启动}\label{ux6bcfux6b21ux4f1aux8bddux542fux52a8}

\begin{enumerate}
\def\labelenumi{\arabic{enumi}.}
\tightlist
\item
  读本文件（你正在做）
\item
  扫一眼 \texttt{knowledge/07\_任务/里程碑.md} 了解当前进度
\item
  扫一眼 \texttt{knowledge/07\_任务/README.md} 了解当前优先级
\end{enumerate}

\subsubsection{执行任务前}\label{ux6267ux884cux4efbux52a1ux524d}

\begin{enumerate}
\def\labelenumi{\arabic{enumi}.}
\tightlist
\item
  先判断是否复杂到需要 plan mode
\item
  能用 agent 并行的事不要串行
\item
  改代码前先跑 \texttt{python\ -m\ pytest\ tests/\ -v\ -\/-tb=short}
  确认基线
\item
  改完后再次跑测试确认不破坏
\end{enumerate}

\subsubsection{理论工作}\label{ux7406ux8bbaux5de5ux4f5c}

\begin{itemize}
\tightlist
\item
  证明优先于实验
\item
  假设必须显式陈述
\item
  审稿人视角前置：每提出一个 claim，先问''审稿人会用什么理由 reject''
\end{itemize}

\subsubsection{论文工作}\label{ux8bbaux6587ux5de5ux4f5c}

\begin{itemize}
\tightlist
\item
  Paper 2 是当前焦点（最接近投稿）
\item
  不要混淆论文的贡献边界（EGP gauge fixing ≠ SCX noise detection）
\item
  公式浓度：Paper 1 低（Nature 风格），Paper 3+4 高（JMLR 风格）
\end{itemize}

\subsubsection{商业化意识}\label{ux5546ux4e1aux5316ux610fux8bc6}

\begin{itemize}
\tightlist
\item
  区分''论文里写什么''和''代码里开源什么''
\item
  L1-L2 全公开，L4-L5 闭源保留
\item
  期刊投稿前确认专利状态
\end{itemize}

\begin{center}\rule{0.5\linewidth}{0.5pt}\end{center}

\subsection{当前阻塞与下一步}\label{ux5f53ux524dux963bux585eux4e0eux4e0bux4e00ux6b65}

\textbf{阻塞}：GPU 未到 → AlN v3 去噪重训、CIFAR 完整训练、MedMNIST 强
backbone 无法进行

\textbf{当前可推进}（无 GPU 依赖）： 1. 🔴 Paper 2 收尾投稿（1-2 周） 2.
🟡 Paper 1 理论部分草稿（定理已就绪） 3. 🟡 UCI tabular 实验（CPU 可做）
4. 🟢 公开 MLIP 数据下载分析（CPU 可做）

\textbf{参考}：\texttt{knowledge/07\_任务/} 下的所有任务文件

\begin{center}\rule{0.5\linewidth}{0.5pt}\end{center}

\emph{本文件是 AI 恢复入口。用户每次与新的 AI 会话时，指示 AI
先读这个文件。}

\end{document}
